\documentclass[11pt]{article}
\usepackage{scriptstyle}
\scripttopic{The Architecture Page}

\begin{document}

\begin{center}
{\LARGE\bfseries CollusionGraph --- Presentation Script}\\[4pt]
{\large The Architecture Page}\\[8pt]
{\small Total running time: \textbf{about 16 minutes}\quad\textbar\quad
19 slides\quad\textbar\quad Runs \emph{after} the dashboard script.}
\end{center}

\vspace{6pt}
\begin{tcolorbox}[colback=metagrey!6,colframe=metagrey!45,boxrule=0.6pt,arc=1pt]
\textbf{Set-up.} Open \texttt{docs/architecture.html} in any browser. It needs
no internet and no server. Use the \textbf{light} theme --- it projects better.
Scroll to the top of the system flowchart before you start.

Same three boxes as the dashboard script. ``\bt'' is a breath.
\stage{square brackets} are things you do, not words you say.

\textbf{The one sentence this whole talk hangs on:} the top half of the diagram
\emph{builds} results and then stops running; the bottom half can only
\emph{read} them.
\end{tcolorbox}

% ===========================================================================
\phase{Part 1 --- Opening (1 minute 15)}

\slide{Slide 1 --- Handing Over --- From The Product To The Machine}{Medium (6/10)}{$\approx$ 35 seconds}

\begin{saythis}
Thanks \blank. I'm \blank, and I'll show you what's underneath that.

And I want to start with something that sounds like a mistake.

Everything you just watched \bt every case, every score, every explanation \bt
was produced by a program that \emph{is not running}. It ran, it wrote its
answers into files, and it shut itself down. What's serving that website can't
train anything and can't change a single score.

That sounds like a limitation. It's the opposite \bt it's deliberate, and it's
the decision I'd defend hardest on this page.

So: this is one page, it reads top to bottom like a production line, and
everything on that console comes out of the bottom of it.

I'll take you down the pipeline first. Then I'll stop at one line in the
middle. And then I'll show you what this looks like hosted properly as a real
product, with the bill.
\end{saythis}

\begin{means}
A handover slide with a hook, so the second speaker doesn't open on
housekeeping.

\textbf{Why the hook works.} ``The thing that made all this isn't running''
sounds wrong for about two seconds, and those two seconds are attention you
would otherwise have spent on ``now I'll explain the architecture''. You then
resolve it immediately, and the room has already been introduced to the trust
boundary before you name it.

Naming that line as ``the decision I'd defend hardest'' makes them wait for it,
which means they are listening when it arrives on slide 10.
\end{means}

% ---------------------------------------------------------------------------
\slide{Slide 2 --- Two Boxes And A Line}{HIGH (8/10)}{$\approx$ 40 seconds}

\begin{saythis}
Before any detail, look at the overall shape. Two big boxes and a line between
them.

The green box at the top is the deep-learning pipeline. It runs, it does its
work, it writes its answers into files, and then it shuts down. It is not
running right now.

The blue box underneath is the live product \bt the website, the service behind
it, the assistant. That part is always on, and it can only \emph{read} what the
green part wrote.

Between them, that dashed line labelled trust boundary. Hold that thought,
because I'm coming back to it.
\end{saythis}

\begin{means}
Get the room to see the shape before any of the boxes. If they leave with only
one thing from this section, it should be ``it trains, then stops, and the
website can only read.''

Everything else on the page --- the security story, the cost story, the fact
that the whole thing demonstrates on a laptop with the wifi off --- follows
from that split.
\end{means}

% ===========================================================================
\phase{Part 2 --- Down the pipeline (5 minutes)}

\slide{Slide 3 --- Boxes 1 and 2 --- Making Two Worlds Look The Same}{HIGH (8/10)}{$\approx$ 1 minute}

\begin{saythis}
Box one, the readers. Six datasets, six readers, one each. And each reader
checks the file's fingerprint before it touches anything. If a download was
cut short, or somebody swapped a file, it stops rather than quietly training on
the wrong data.

Box two is the one I'd underline on the whole page.

Everything gets converted into one common format. Things, connections between
things, and where we have them, answers. And the time-splitting rules get
applied right here.

Here's why box two is the trick of the entire project. After that box, the
system genuinely doesn't know which crime it's looking at. A bank account and a
construction company arrive as the same kind of object.

That's what lets one system handle both worlds. And it's the only reason we
could even \emph{ask} whether knowledge of one crime transfers to the other.
\end{saythis}

\begin{means}
\textbf{The everyday version.} Two countries with different plug sockets. You
could build two of every appliance --- or you could build one adapter. Box two
is the adapter, and everything downstream is built once instead of twice.

\textbf{The research point.} You cannot ask ``does money-laundering knowledge
transfer to bid rigging?'' unless both are written in the same language. The
shared format is what makes the question askable. Say that --- it turns a
plumbing box into a research contribution.

\textbf{The fingerprint check} is called a checksum. It is one line of code and
it is the difference between ``our numbers are wrong'' and ``the pipeline
refused to start''.
\end{means}

\begin{asked}
\q{Don't you lose information flattening two domains into one format?} You lose
nothing you cared about --- domain-specific facts survive as extra attributes.
What you gain is that everything after that point is written once and audited
once, instead of twice with two sets of bugs.
\end{asked}

% ---------------------------------------------------------------------------
\slide{Slide 4 --- Box 3 --- Four Opinions At The Same Time}{VERY HIGH (9/10)}{$\approx$ 1 minute 20}

\begin{saythis}
Box three is four boxes side by side, and they're side by side because they
genuinely run at the same time on the same network. Four different opinions
about the same accounts.

The first group learns from examples where somebody already knew the answer.
Three designs, and each one cares about something different. One works
out which of your connections matter most and listens hardest to those. One
treats money coming in differently from money going out. One handles different
\emph{kinds} of relationship separately \bt winning a contract isn't the same
as buying from someone.

The second group needs no answers at all. It learns what normal looks like, and
flags whatever it can't rebuild. Hold on to that one, because I'll come back to
it and it's the most interesting result we have.

The third group is hand-computed measurements. How many connections does this
thing have. What shapes is it sitting inside. And for contracts, the classic
warning signs \bt the same firms always bidding together, prices sitting
suspiciously close.

And the fourth is the honesty group. The strongest old-fashioned method we
could field, plus the planted-cartel tests, plus the automated test suite. It's
there so nobody can accuse us of only racing against ourselves.
\end{saythis}

\begin{means}
Four arms, four different failure modes, and they don't overlap. That is the
argument for having four.

\textbf{Learning without answers, in one sentence.} It studies the ordinary
traffic until it can reproduce it from memory; anything it can't reproduce is
unusual. No answer key required.

\textbf{Why the honesty arm matters to a marker.} A deep-learning result with
no strong classical comparison is a demonstration, not a result. Arm four is
also why you're able to admit, later and without flinching, that on one dataset
the old-fashioned method wins.
\end{means}

\begin{asked}
\q{Why keep the no-answers arm if it scores worst?} Because it's the only arm
that works where there is no answer key --- which is the realistic deployment
case. I'll show you the number in two slides, and it's the strongest
deep-learning result in the project.
\end{asked}

% ---------------------------------------------------------------------------
\slide{Slide 5 --- The No-Answers Arm --- Our Best Result}{VERY HIGH (9/10)}{$\approx$ 1 minute}

\begin{saythis}
Let me pay off that promise now, because it's the experiment I'd want you to
remember.

We took a real national contracts database \bt about a hundred and sixty
thousand entities \bt and we planted a hundred fake cartels inside it. Nine
hundred and forty companies, five different ways of cheating. And critically,
we gave the detectors \emph{no answer key at all}.

Then we asked: reviewing about one in eighty of that network, how many do you
find?

The result splits cleanly, and both halves are useful. Cartels that physically
bid together \bt so they leave a visible web of connections \bt we find about
nine times in ten. Cartels that just take turns winning, and never appear
together anywhere, escape almost completely.

And here's the part that matters for this project. In that experiment the
old-fashioned method can't even take part, because there's nothing to train it
on. The simple non-learning baseline scores about one in eight. The deep model
scores about nine in ten.

That's deep learning doing something nothing else in the toolbox can do.
\end{saythis}

\begin{means}
This is your answer to the single hardest question in the project --- ``your
XGBoost beats your neural network, so why bother?'' --- and it is much stronger
delivered proactively here than defensively later.

\textbf{Why knowing what it \emph{misses} is a feature.} A deployment team
would far rather hear ``this catches bid-together cartels and is blind to
take-turns cartels'' than ``it is eighty-two percent accurate''. The first
tells them what to do next; the second tells them nothing.

\textbf{The honest framing to use.} If the cheats leave a visible web between
them, we find them. If they simply agree to take turns and never appear
together, shape alone cannot see it --- you would need prices and dates.
\end{means}

\begin{asked}
\q{Aren't planted cartels easier than real ones?} They're generated from the
published descriptions of how these cartels actually operate, and three of the
five kinds escape us completely --- which is not what you'd see if we'd made
them easy.

\q{Why not just use real labelled cartels?} There are almost none. Confirmed
cartels are rare, and that scarcity is the whole reason the no-answers arm has
to exist.
\end{asked}

% ---------------------------------------------------------------------------
\slide{Slide 6 --- Box 4 --- Why The Word ``First'' Is Doing All The Work}{VERY HIGH (9/10)}{$\approx$ 1 minute}

\begin{saythis}
Box four merges those four opinions into one number per account. And the word
``first'' in that box is doing enormous work.

Here's the problem. Two models can both be excellent at ranking and still be on
completely different scales. One says nought point nine when it's fairly
confident. The other says nought point one when it's absolutely certain.
Average those raw, and the loud one flattens the quiet one, regardless of which
one was right.

So we translate each model onto a proper scale first, checked against held-out
data, and only then combine them.

And we didn't take that on faith \bt we ran it both ways. Skip the translation
step and the combined result collapses to roughly what blind guessing gets you.
Do it properly and it's about nine times better.

Same models. Same data. That single decision about \emph{ordering} was worth
more than any model we trained.
\end{saythis}

\begin{means}
\textbf{The everyday version.} Three judges scoring a competition. One marks
everything between six and seven. One uses the full range. Average their raw
scores and the cautious judge is silently ignored. You normalise them first.

\textbf{Why this slide is worth a full minute.} It is a genuine measured
finding, it is easy to understand, and it is the kind of engineering detail
that separates a project that ran experiments from a project that ran a
tutorial.

The technical name is \emph{calibration}. Say it only if asked.
\end{means}

\begin{asked}
\q{Why combine at all instead of using the best single model?} On one dataset
we do report the best single model, and it's the old-fashioned one. The
combination is what protects you when one member is weak --- and we measured
that combining badly is far worse than not combining at all.
\end{asked}

% ---------------------------------------------------------------------------
\slide{Slide 7 --- Box 5 --- Turning Scores Into Somebody's Morning}{HIGH (8/10)}{$\approx$ 50 seconds}

\begin{saythis}
Box five is where research turns into a product.

Up to here we've got a score on every single account. Sixty-seven thousand of
them on the money side. That's useless to a human being.

So we find the clusters. We rank them. We throw away near-duplicates so nobody
opens the same case twice. And we cut the list at whatever number of cases the
team can actually review.

Two hundred and fifty-four cases on the money side. Two hundred and twenty-three
on the contracts side. That's a morning's work, not a research output.

And notice what this step does \emph{not} use. The answers. It runs on
structure alone \bt which is exactly why the identical step also runs on the
data where no answers exist.
\end{saythis}

\begin{means}
This box is the bridge back to the console the room already saw, so gesture at
the numbers they recognise --- 254 and 223 were on the Overview screen.

The ``it never uses the answers'' point is what makes the planted-cartel study
possible at all, so it is worth ten seconds even though it sounds like a
technicality.
\end{means}

% ---------------------------------------------------------------------------
\slide{Slide 8 --- Box 6 --- Why An Alert Alone Isn't Enough}{HIGH (8/10)}{$\approx$ 1 minute}

\begin{saythis}
Box six. Three things happen to every case at the top of the list.

First, the explainer picks out the handful of connections the decision actually
depended on \bt out of a group of eighty-five accounts, maybe three of them and
four payments. And then we check that, by re-running the case with only those
connections, and again with them taken away.

Second, a set of hand-written rules checks the case against nine known shapes
and puts a name on it \bt but only when it can actually prove the shape.

Third, when a shape matches, we attach the official warning sign it corresponds
to. Quoted, with the source stamped on it.

And those three are deliberately independent. The explainer is learned and can
be wrong. The rules are fixed and cannot invent anything. So an investigator
who distrusts the neural network entirely can throw the first one away, and the
other two still stand.
\end{saythis}

\begin{means}
\textbf{Define ``the handful of connections'' before you use it.} That phrase
was a real failure in our own report --- it appeared before it was ever
defined, so the whole explanation collapsed for the reader. The definition is:
out of everything in the group, the specific few connections the system says
its decision depended on.

\textbf{The everyday version of the check.} Someone says two exam questions
made them fail. Show them only those two --- do they still fail? Delete those
two --- do they now pass? Both yes means those questions really were the
reason.

The numbers, if pushed: this check used to fail on thirty-eight of our fifty
published explanations. We replaced the explanation method and now it fails on
one.
\end{means}

\begin{asked}
\q{How do you know the explanation is honest and not a story?} Because it's
tested rather than asserted, by those two re-runs. And when the two answers
contradict each other, we put a warning on the case file instead of hiding it.
\end{asked}

% ---------------------------------------------------------------------------
\slide{Slide 9 --- Box 7 --- Just Files}{Medium (5/10)}{$\approx$ 30 seconds}

\begin{saythis}
Box seven. Everything the pipeline produced goes into ordinary files in a
folder. Not a database. Not a running service. Files.

That sounds unambitious, and it's one of the better decisions in the project.
Files can be copied, fingerprinted, versioned, handed to an examiner on a
memory stick, and undone by pointing at a different folder.

And a file opened read-only cannot be changed by the thing reading it \bt
which is exactly what the next slide is about.
\end{saythis}

\begin{means}
This is the set-up for the trust boundary. Deliver the last line and pause
before moving on.
\end{means}

% ===========================================================================
\phase{Part 3 --- The most important line on the page (1 minute 15)}

\slide{Slide 10 --- The Trust Boundary}{VERY HIGH (10/10)}{$\approx$ 1 minute 15}

\begin{saythis}
Here it is. This is the line I said I'd come back to.

Above it: the part that trains. That part is powerful, because it can change
what the system believes about a person. Below it: the part the outside world
can actually reach. That part is attackable, because anything on the internet
is.

Powerful and attackable must never be the same process. So we cut them
completely apart. The top half writes files and exits. The bottom half opens
those files read-only. There is no route back upwards.

Now think about what that means if somebody breaks into our website tomorrow.

They can read published results. They cannot retrain a model. They cannot
change anybody's score. They cannot plant evidence in a case file.

And that's the attack that would actually matter here. Not stealing the data
\bt quietly changing who looks suspicious. That's much harder to notice and far
worse in its consequences.

There's a bonus too. Because the serving side never trains anything, it needs
no expensive hardware. It starts instantly, runs on the cheapest machine on the
price list, and demonstrates on a laptop with the wifi switched off.

The security argument and the cost argument turn out to be the same argument.
\end{saythis}

\begin{means}
This is the highest-value slide in the architecture talk. Rehearse it until it
is smooth, and do not rush the ``quietly changing who looks suspicious'' line
--- that is the sentence that shows you thought about the domain and not just
the code.

\textbf{The strongest version of the claim,} if someone pushes: in the cloud
design this is enforced by the cloud provider's own permission system, not by
our code being well-behaved. Our code could contain a bug that tries to write,
and the write would still be refused.
\end{means}

\begin{asked}
\q{Isn't that just a deployment detail?} In a system whose output decides who
gets investigated, ``an attacker cannot alter a score'' isn't a deployment
detail, it's an integrity property. It's also why the service has no write
function at all --- there's nothing to abuse, because none exists.

\q{What if the attacker gets into the training side instead?} Then they have a
problem worth having, because that side isn't running. It exists for a couple
of hours when a scheduled job starts it, and it isn't reachable from the
internet at all.
\end{asked}

% ===========================================================================
\phase{Part 4 --- Below the line (1 minute 20)}

\slide{Slide 11 --- Box 8 --- The Doorway, And Where The Assistant Sits}{HIGH (7/10)}{$\approx$ 50 seconds}

\begin{saythis}
Below the line, box eight. A small read-only service that hands out those
files. Ask it for a case, it gives you the case. Ask it for a network, and it
trims that network down to a sensible size \emph{on the server} before sending,
so your browser never receives a million connections.

Every single response carries the ethics wording. Not just the pretty pages
\bt the raw data too.

And beside it, in its own box, the assistant. Notice it sits \emph{beside} the
doorway rather than behind it, and both arrows point in and out of that same
read-only service. It has no private access to anything. It reads exactly what
a person looking at the dashboard could read.
\end{saythis}

\begin{means}
The placement on the diagram is the argument. If the assistant had its own line
into the data store, everything you said about guardrails would rest on trust.
Because it goes through the same read-only doorway, it rests on architecture.
\end{means}

% ---------------------------------------------------------------------------
\slide{Slide 12 --- Boxes 9 and 10 --- The Console, And The Person}{HIGH (7/10)}{$\approx$ 30 seconds}

\begin{saythis}
Box nine is the console you've already seen.

Box ten is a person. And it's drawn on the architecture diagram on purpose.

The system produces a ranked list and the evidence behind it. A human decides
what happens next, inside a fixed number of cases a week.

That's not modesty about how good the model is. If a screening tool made the
escalation decision, you'd have automated an accusation.
\end{saythis}

\begin{means}
Putting a human on an architecture diagram is unusual, and that is exactly why
it works. It makes the human-in-the-loop constraint structural rather than
aspirational.

``You'd have automated an accusation'' is the line to land. Then move on.
\end{means}

% ===========================================================================
\phase{Part 5 --- The method cards (1 minute 30)}

\slide{Slide 13 --- The Cards Below The Flowchart}{Medium (6/10)}{$\approx$ 1 minute 30}

\begin{saythis}
Scroll past the flowchart and every box gets a card explaining what it does and
how it works. I won't read twelve cards at you. Three are worth pausing on.

The graph networks. Ordinary machine learning looks at a thing's own details.
These look at \emph{who it deals with}. Each account starts with a profile of
numbers, and then in a few rounds every account blends in its neighbours'
profiles. After a couple of rounds, an account sitting inside a laundering ring
looks measurably different from a normal one \bt even if its own details looked
perfectly ordinary.

The rare-crime problem. Only about one in fifty of these things is a real case.
A lazy model calls everything innocent, is right ninety-eight percent of the
time, and is completely useless. So we change what the training rewards. Easy
correct answers earn almost nothing, and the effort goes where the rare cases
are.

And the old-fashioned yardstick. Hundreds of small decision trees, each
correcting the last. On the Bitcoin data they still beat our neural networks,
and we publish that with a proper statistical test rather than burying it.
\end{saythis}

\begin{means}
\textbf{The everyday version of a graph network.} Deciding whether a new pupil
is trouble. You could read their file. Or you could notice who they eat lunch
with every day. After a few days of watching, the lunch table tells you more
than the file.

\textbf{The rare-crime fix} is called focal loss. Only name it if asked.

\textbf{Volunteering that the trees win} is a strength. It reads as confidence,
and it sets up the answer you already gave on slide 5 --- the trees cannot
enter the experiment that matters most.
\end{means}

\begin{asked}
\q{The trees beat your deep models --- doesn't that sink the project?} On that
one dataset, yes, and we say so. That dataset already ships hand-built
neighbour summaries as ordinary columns, so the trees get the network
information for free. Take the answer key away, as in the planted-cartel study,
and they can't run at all.
\end{asked}

% ===========================================================================
\phase{Part 6 --- The cloud, and the bill (3 minutes 30)}

\slide{Slide 14 --- The Always-On Half}{HIGH (8/10)}{$\approx$ 1 minute 10}

\begin{saythis}
Second diagram. What this looks like hosted properly. Same two halves, same
line between them.

Top half is always on. Users need a browser and nothing else.

Three things happen before anything of ours runs. An address service turns our
web address into a real machine. A delivery network keeps copies of the site
around the world, so it loads from somewhere near you rather than from one
machine on one continent. And a filter in front of everything blocks the
well-known attacks.

The dashboard itself is just files in cheap storage, in a private area only the
delivery network is allowed to read.

The results live in their own storage, and this is the box I'd point at. Every
run writes into its own dated folder. Nothing is ever overwritten.

Which means undoing a bad model is pointing at yesterday's folder. Seconds. No
retraining, no panic \bt and the bad run is still sitting there to be
investigated rather than destroyed.
\end{saythis}

\begin{means}
The versioned-folder point is the one an examiner with industry experience will
respond to, because rollback is the operational problem everybody has actually
lived through.

\textbf{Why the site bucket is private.} A public bucket is one
misconfiguration away from being writable by somebody else, and it lets people
bypass the protective filter. Making the delivery network the only reader means
every visitor goes through the filter without exception.
\end{means}

\begin{asked}
\q{Why not use a proper database for the results?} Because we don't need one.
The results are written once and read many times, never edited. Files give us
versioning, checksums and rollback for free --- and a file opened read-only
cannot be modified by the website reading it, which is the trust boundary
again.
\end{asked}

% ---------------------------------------------------------------------------
\slide{Slide 15 --- The Pieces That Exist So Nothing Goes Quietly Wrong}{Medium (6/10)}{$\approx$ 50 seconds}

\begin{saythis}
Along the bottom of that box, four small pieces that aren't glamorous and all
of them earn their place.

The assistant's key lives in a secrets store, so it's never in the code and
never in the repository.

The public side is granted \emph{read-only} rights on the results. And that's
the strongest sentence in this diagram \bt the boundary I showed you isn't just
our own code being polite. Amazon enforces it. Our code could have a bug that
tries to write, and the write would still be refused.

Logs and alarms, so you find out something broke before a user tells you.

And a spending alarm, switched on from day one, so a stuck job can't quietly
cost a fortune while nobody's looking.
\end{saythis}

\begin{means}
Say the read-only permissions line slowly. It converts a claim into a
guarantee, and it is the single most defensible thing on the cloud diagram.

The spending alarm sounds like a joke until you mention it is a student
project. Then it sounds like judgement.
\end{means}

% ---------------------------------------------------------------------------
\slide{Slide 16 --- The Half That Only Exists While It Runs}{HIGH (8/10)}{$\approx$ 55 seconds}

\begin{saythis}
Bottom half, and this is the money-saving idea.

There's an alarm clock. Every night, or every week, it starts the pipeline.
That job rents a graphics computer at about seventy percent off, reads the
data, trains, scores, groups, explains, writes the results \bt and then shuts
itself down.

Now, renting at seventy percent off means the machine can be taken away from
you with almost no warning. That would be unacceptable for anything serving
live traffic. But this job serves nothing. If it gets interrupted, the website
carries on showing yesterday's results, an alarm fires, and it runs again
tonight.

The fact that this job can be interrupted safely is exactly what earns the
discount.

And follow that arrow going up. It's the only connection between the two
halves. The nightly job \emph{writes} into the storage the website
\emph{reads}. There's nothing coming back down.
\end{saythis}

\begin{means}
The spot-pricing justification is the good bit. Anyone can say ``we use cheap
machines''. Explaining \emph{why this workload is the one that can tolerate
cheap machines} is engineering judgement.

The upward arrow is the trust boundary again, drawn in cloud services. Point at
it and say so --- repetition of the central idea across two diagrams is a
feature.
\end{means}

\begin{asked}
\q{What if it gets taken away halfway through?} The website keeps serving
yesterday's results, because nothing was overwritten in place. A half-finished
run just produces an incomplete folder that nobody is pointing at.
\end{asked}

% ---------------------------------------------------------------------------
\slide{Slide 17 --- What It Costs}{HIGH (7/10)}{$\approx$ 55 seconds}

\begin{saythis}
Three levels, real list prices.

A university demo: website hosting, one small always-on machine, a web address,
and the training done on your own computer or free Colab. About ten dollars a
month \bt and effectively nothing, because standard new-account credits cover
it.

A real launch: a pay-per-use doorway good for about a million requests,
storage, a rented graphics machine roughly four hours a week, pay-per-use
assistant, monitoring. Somewhere between about fifteen and fifty dollars a
month.

A growing product: add containers across two data centres, a load balancer, a
firewall, a nightly graphics machine, team-scale assistant usage. Roughly
ninety to a hundred and eighty.

And notice what actually changes between those three. The machine learning
doesn't. Same pipeline, same results, every level. What you're buying as you go
up is tolerance for things going wrong.
\end{saythis}

\begin{means}
That last paragraph is the insight, and most groups miss it. The tiers are not
``more model'' --- they are more resilience. Tier one assumes somebody will
notice if it stops. Tier three survives a data centre failure.

\textbf{If asked what the risky line is:} the assistant, because it is charged
per unit of text and driven by humans rather than by us. The graphics machine
is the biggest cost but it is bounded --- it runs for a few hours and stops.
\end{means}

\begin{asked}
\q{Could you actually afford to run this?} Yes, and that's a consequence of the
architecture rather than luck. The expensive hardware exists for two hours a
night and then disappears. Showing results needs none of it.
\end{asked}

% ===========================================================================
\phase{Part 7 --- Close (35 seconds)}

\slide{Slide 18 --- The Close}{HIGH (8/10)}{$\approx$ 35 seconds}

\begin{saythis}
So: one page, two halves, one line between them.

Above the line we train, and then we stop. Below the line we show people
things, and we can't change anything. That single decision gave us the security
property, the cost profile, and the ability to demonstrate this on a laptop
with no internet \bt all three from the same choice.

And at the very bottom of the diagram, a person still decides.
\end{saythis}

\begin{means}
Mirror the dashboard script's close. Two talks, same final note, delivered by
different people --- that reads as a group that actually planned together.
\end{means}

% ---------------------------------------------------------------------------
\slide{Slide 19 --- Appendix --- The Four Minute Version}{---}{$\approx$ 4 minutes}

If you are cut short, run slides \textbf{2, 4, 5, 10, 16, 18}.

That is: the shape, the four arms, the result that only deep learning can
produce, the trust boundary, the cloud version of the boundary, and the close.

\vspace{8pt}
\textbf{Never read a decimal aloud on this page:}

\begin{center}\small
\begin{tabular}{@{}p{6.2cm}p{8.4cm}@{}}
\hline
\textbf{On the page} & \textbf{Say} \\
\hline
Combined 0.4434 vs raw 0.0511 & ``about nine times better, same models, same
data'' \\
Guessing level 0.065 & ``blind guessing gets about six in a hundred'' \\
About 2\% positives & ``about one in fifty is a real case'' \\
Planted study, 0.92 vs 0.12 & ``about nine in ten, against about one in eight
for the simple baseline'' \\
Explanation check 38/50 $\to$ 1/50 & ``it used to fail on most of them, now it
fails on one'' \\
24 of 24 exam questions & ``a twenty-four question exam, and all of them have
to pass'' \\
Tier 1 about \$10 & ``about ten dollars a month'' \\
Spot discount $\approx$ 70\% & ``about seventy percent off'' \\
\hline
\end{tabular}
\end{center}

\end{document}
